% KOMA-Script Artikel-Praeambel passend zur AVVP-Beamer-Vorlage (Magazin-Look)
% Engine: LuaLaTeX (wegen fontspec und lokaler OTF/TTF)

\usepackage[ngerman]{babel}
\usepackage{fontspec}
\usepackage{xcolor}
\usepackage{graphicx}
\usepackage{csquotes}
\usepackage{microtype}
\usepackage{hyperref}

% Bibliographie wie im Beamer-Projekt
\usepackage[style=authoryear]{biblatex}
\addbibresource{bibtex.bib}

\usepackage{booktabs}
\usepackage{etoolbox}
\usepackage{float}

% Optional: Seitenlayout
\usepackage[a4paper,margin=2.5cm]{geometry}

% Kopf-/Fußzeile (Magazin)
\usepackage[automark,headsepline]{scrlayer-scrpage}
\usepackage{lastpage}

% Kopf-/Fußzeilen-Inhalt
\clearpairofpagestyles
\automark{section}
\ihead{\headmark}

% Rechts oben: Titel (Kurztitel, falls via \title[Kurztitel]{...} gesetzt)
\makeatletter
\newcommand{\AVVPHeaderTitle}{%
  \ifcsname shorttitle\endcsname
    \ifx\shorttitle\@empty
      \@title
    \else
      \shorttitle
    \fi
  \else
    \@title
  \fi
}
\makeatother
\ohead{\AVVPHeaderTitle}

\cfoot{Seite\ \pagemark\ von\ \pageref{LastPage}}
\pagestyle{scrheadings}

% Farben (aus dem Beamer-Theme uebernommen)
\definecolor{AVVPBg}{RGB}{5,23,41}
\definecolor{AVVPBlue}{RGB}{46,108,198}
\definecolor{AVVPSpark}{RGB}{235,156,70}
\definecolor{AVVPGrey}{RGB}{90,98,110}

% Schriften (lokal aus dem Projekt)
\setmainfont[
  Path = fonts/exo-2/,
  Extension = .otf,
  UprightFont = *-Regular,
  ItalicFont  = *-Italic,
  BoldFont    = *-Bold,
  BoldItalicFont = *-BoldItalic,
  Ligatures = TeX
]{Exo2}

\setsansfont[
  Path = fonts/exo-2/,
  Extension = .otf,
  UprightFont = *-Regular,
  ItalicFont  = *-Italic,
  BoldFont    = *-Bold,
  BoldItalicFont = *-BoldItalic,
  Ligatures = TeX
]{Exo2}

\newfontfamily\AVVPTitleFont[
  Path = fonts/share_techmono/,
  Extension = .otf,
  UprightFont = Share-TechMono,
  Ligatures = {TeX,NoCommon}
]{Share Tech Mono}

% Ueberschriften in CI-Farbe
\addtokomafont{disposition}{\color{AVVPBlue}\AVVPTitleFont}

% AVVP Link-Farbe
\hypersetup{colorlinks=true,linkcolor=AVVPBlue,urlcolor=AVVPBlue,citecolor=AVVPBlue}

% Default fuer \includegraphics: volle Textbreite, Seitenverhaeltnis behalten
\setkeys{Gin}{width=\linewidth,keepaspectratio}



% --- Magazin-Look: weniger Weißraum + ruhigere Captions ---
\setlength{\textfloatsep}{6pt plus 2pt minus 2pt}
\setlength{\intextsep}{6pt plus 2pt minus 2pt}
\setlength{\floatsep}{6pt plus 2pt minus 2pt}
\setlength{\abovecaptionskip}{3pt}
\setlength{\belowcaptionskip}{0pt}

% Abbildungen/Tabellen im Artikel möglichst "hier" setzen
\floatplacement{figure}{H}
\floatplacement{table}{H}

% Linksbuendig in figure/table, auch wenn im Dokument \centering steht
\AtBeginEnvironment{figure}{\let\centering\raggedright}
\AtBeginEnvironment{table}{\let\centering\raggedright}

% Captions kleiner, grauer und linksbuendig
\makeatletter
\long\def\@makecaption#1#2{%
  \vskip\abovecaptionskip
  {\footnotesize\color{AVVPGrey}#1: #2\par}%
  \vskip\belowcaptionskip
}
\makeatother

% ------------------------------------------------------------
% Bildnachweis (Artikel): kompatible Kurzform wie im Beamer (cite-key)
%   Nutzung: \AVVPCreditInline{cite-key}
% ------------------------------------------------------------
\makeatletter
\newcommand{\AVVPCreditInline}[1]{%
  \begingroup
    \@ifpackageloaded{biblatex}{\nocite{#1}}{}%
    \;\textnormal{([#1])}%
  \endgroup
}
\makeatother

% ------------------------------------------------------------
% VIDEO HELPERS (wie in der Beamer-Präambel; Artikel-Implementierung)
%   Ziel: gleiche Makro-Namen/Signaturen, aber Magazin-taugliche Ausgabe.
% ------------------------------------------------------------

% Poster-Image (skalieren ohne Verzerrung, Hoehe begrenzen)
\newcommand{\AVVPVideoPoster}[2][]{%
  \includegraphics[
    width=0.5\linewidth,
    height=0.35\textheight,
    keepaspectratio,
    #1%
  ]{#2}%
}

% Minimaler File-URI Builder (ersetzt Leerzeichen durch %20)
\makeatletter
\ExplSyntaxOn
\tl_new:N \g_avvp_file_uri_tl
\cs_new_protected:Npn \avvp_build_file_uri:n #1
  {
    \tl_set:Nn \l_tmpa_tl {#1}
    \tl_replace_all:Nnn \l_tmpa_tl {\c_space_tl} {\c_percent_str 20}
    \tl_gset:Nx \g_avvp_file_uri_tl { run:\tl_use:N \l_tmpa_tl }
  }
\newcommand{\AVVPBuildFileURI}[1]{%
  \begingroup
    \avvp_build_file_uri:n {#1}%
    \xdef\AVVPFileURI{\tl_use:N \g_avvp_file_uri_tl}%
  \endgroup
}
\ExplSyntaxOff
\makeatother

% Label wie in Beamer: nur "ordner/datei" (fallback: kompletter String)
\makeatletter
\ExplSyntaxOn
\cs_new_protected:Npn \avvp_video_local_label:n #1
  {
    \seq_set_split:Nnn \l_tmpa_seq {/} {#1}
    \int_compare:nNnTF {\seq_count:N \l_tmpa_seq} > {1}
      {
        \tl_set:Nx \l_tmpb_tl {\seq_item:Nn \l_tmpa_seq {-2}}
        \tl_set:Nx \l_tmpc_tl {\seq_item:Nn \l_tmpa_seq {-1}}
        \texttt{\nolinkurl{\tl_use:N \l_tmpb_tl/\tl_use:N \l_tmpc_tl}}
      }
      {\texttt{\nolinkurl{#1}}}
  }
\newcommand{\AVVPVideoLocalLabel}[1]{\avvp_video_local_label:n {#1}}
\ExplSyntaxOff
\makeatother

% Clickable poster -> URL
\newcommand{\AVVPVideoURL}[3][]{%
  \begin{center}%
    \href{#2}{\AVVPVideoPoster[#1]{#3}}%
  \end{center}%
}

% Clickable poster -> lokale Datei
\newcommand{\AVVPVideoLocal}[3][]{%
  \begin{center}%
    \AVVPBuildFileURI{#2}% setzt \AVVPFileURI
    \href{\AVVPFileURI}{\AVVPVideoPoster[#1]{#3}}\\[0.4ex]%
    {\scriptsize\href{\AVVPFileURI}{\AVVPVideoLocalLabel{#2}}}%
  \end{center}%
}

% ------------------------------------------------------------
% Quellenverzeichnis wie im Beamer: cite-key [..] am Anfang jeder Zeile
% ------------------------------------------------------------
\makeatletter
\AtBeginDocument{%
  \@ifpackageloaded{biblatex}{%
    \DeclareFieldFormat{entrykey}{[#1]}%
    \renewbibmacro*{begentry}{\printfield{entrykey}\addspace}%
  }{}%
}
\makeatother

% Bildpfade wie in der Beamer-Vorlage (Aufruf aus AVVP_Beamer/)
\graphicspath{{./}{pics/}{media/}}

% ------------------------------------------------------------
% AVVP Titel-Assets (Benennung analog zur Beamer-Praeambel)
% ------------------------------------------------------------
\def\AVVPTitleLogoDarkFile{pics/avvp_2019_logo_wortmarke_neg.png}
\def\AVVPTitleLogoLightFile{pics/avvp_2019_logo_wortmarke_pos_gross.png}
\def\AVVPTitleSloganText{Faszination. Astronomie. Erleben.}

% Modus-Flag analog zum Beamer-Theme
\newif\ifAVVPIsLightMode
% Artikel hat weissen Hintergrund -> default Light
\AVVPIsLightModetrue

\newcommand{\AVVPDarkMode}{\global\AVVPIsLightModefalse}
\newcommand{\AVVPLightMode}{\global\AVVPIsLightModetrue}

\newcommand{\AVVPTitleLogoAuto}{%
  \ifAVVPIsLightMode
    \includegraphics[width=0.30\textwidth]{\AVVPTitleLogoLightFile}%
  \else
    \includegraphics[width=0.30\textwidth]{\AVVPTitleLogoDarkFile}%
  \fi
}

\newcommand{\AVVPTitleSloganColor}{%
  \ifAVVPIsLightMode
    \color{AVVPBg}% dark text on light background
  \else
    \color{white}% light text on dark background
  \fi
}

% Titelseite fuer KOMA-Artikel (Logo + Drei-Zeiler)
\makeatletter
\newcommand{\AVVPMakeTitlePage}{%
  \begin{titlepage}
    \centering

    {\AVVPTitleFont\color{AVVPBlue}\LARGE\@title\par}

    \ifx\@subtitle\@empty\else
      \vspace{0.8em}
      {\AVVPTitleFont\color{AVVPSpark}\large\@subtitle\par}
    \fi

    \vspace{1.2em}
    {\color{AVVPBlue}\begingroup\let\and\\\@author\par\endgroup}

    \vspace{0.6em}
    {\color{AVVPBlue}\@date\par}

    \vfill

    \AVVPTitleLogoAuto\par

    \vspace{0.8em}
    {\AVVPTitleFont\AVVPTitleSloganColor\normalsize\AVVPTitleSloganText\par}

  \end{titlepage}
}
\makeatother